\documentclass[11pt,a4paper]{article}

\usepackage[margin=0.85in]{geometry}
\usepackage{amssymb,amsmath}
\usepackage{enumitem}
\usepackage{booktabs}
\usepackage{parskip}

\setlist[itemize]{leftmargin=*,itemsep=0.15em,topsep=0.25em}
\setlist[enumerate]{leftmargin=*,itemsep=0.15em,topsep=0.25em}

\title{\textbf{Math Bootcamp Syllabus}\\[0.35em]
\large For Incoming MSc Economics Students}
\author{Math Camp 2026}
\date{}

\begin{document}

\maketitle
\vspace{-1.2em}

\section*{Overview}

\begin{itemize}
  \item \textbf{Objective.} Bridge a weaker undergraduate math background to MSc-level economics (micro, macro, econometrics, mathematical economics).
  \item \textbf{Format.} 5 days $\times$ 3 hours. Emphasis on intuition and standard problem templates; only simple proofs where needed. Detailed proofs live in slide appendices / fall courses.
  \item \textbf{Materials.} Day-by-day lecture slides; this syllabus is the concise roadmap.
\end{itemize}

\section*{Day-by-day topics}

\subsection*{Day 1 --- Foundations (\S1.1--1.9)}
\begin{enumerate}
  \item \textbf{Notation, logic, sets, functions.} Number systems; set membership and logic symbols; necessary/sufficient/contrapositive; set operations and Cartesian products; domain/codomain/range; injective/surjective.
  \item \textbf{Sup/inf and algebra.} $\max/\min$ vs.\ $\sup/\inf$; logs, elasticities, Cobb--Douglas; homogeneity and Euler's theorem (CRS adding-up).
  \item \textbf{Calculus.} Single-variable derivatives (rules, marginal interpretation); multivariate: partials, total differential, gradient, directional derivatives, cross-partials; limits, continuity, differentiability.
\end{enumerate}

\subsection*{Day 2 --- More calculus and analysis bridge (\S2.1--2.6)}
\begin{enumerate}
  \item \textbf{Implicit function theorem.} Comparative statics when $y^*(\theta)$ is only defined implicitly; tax incidence.
  \item \textbf{Existence tools and integration.} Open/closed/compact sets; EVT / Weierstrass; MVT; Riemann integral and FTC; expectation via CDF (integration by parts).
  \item \textbf{Leibniz, Fubini, Taylor, uniform convergence.} Differentiating under the integral; changing order of integration; linearization / Taylor; why uniform convergence matters for exchanging limits.
\end{enumerate}

\subsection*{Day 3 --- Linear algebra and real analysis}
\begin{enumerate}
  \item \textbf{Linear algebra.} Vectors, matrices, $Ax=b$, transpose, determinants ($2\times 2$), invertibility, linear independence, rank.
  \item \textbf{Real analysis bridge.} Uniform-convergence practice; separation / supporting hyperplanes (welfare intuition); fixed-point existence statements (Brouwer / contraction, no proofs).
\end{enumerate}

\subsection*{Day 4 --- Optimization (\S4.1--4.5)}
\begin{enumerate}
  \item \textbf{Unconstrained optima.} FOC/SOC; concavity via $f''$ / Hessian; local vs.\ global.
  \item \textbf{Shape and equality constraints.} Concave vs.\ quasiconcave; convex upper contour sets; Lagrange multipliers; shadow prices; envelope theorem.
  \item \textbf{Inequalities and KKT.} Binding vs.\ slack; complementary slackness; geometric Farkas $\Rightarrow$ KKT; corner solutions (nonnegativity).
\end{enumerate}

\subsection*{Day 5 --- Dynamic programming}
\begin{enumerate}
  \item \textbf{Framework.} Sequential decisions; state vs.\ control; law of motion; finite-horizon focus and Markov idea.
  \item \textbf{Bellman and solution method.} Value function and principle of optimality; Bellman equation; backward induction; policy function.
  \item \textbf{Worked core.} Two-period consumption--savings; Euler equation; link to standard intertemporal optimization; VFI intuition (infinite horizon, conceptual only).
\end{enumerate}

\section*{Cross-cutting goals}
\begin{itemize}
  \item Read FOCs, SOCs, and comparative statics fluently.
  \item Map math objects to economic meaning (marginals, shadow prices, value functions).
  \item Build confidence with the templates used in fall core courses.
\end{itemize}

\section*{Suggested time allocation}

\begin{center}
\footnotesize
\setlength{\tabcolsep}{4pt}
\begin{tabular}{@{}clp{0.24\textwidth}p{0.24\textwidth}p{0.24\textwidth}@{}}
\toprule
Day & Theme & Block 1 ($\sim$1\,hr) & Block 2 ($\sim$1\,hr) & Block 3 ($\sim$1\,hr) \\
\midrule
1 & Foundations & Notation, logic, sets, functions & Sup/inf; algebra; Cobb--Douglas & Derivatives; Euler; multivariate; limits \\
2 & Calculus bridge & IFT & EVT/MVT; integration & Leibniz/Fubini; Taylor; uniform conv. \\
3 & Lin.\ alg.\ /\ analysis & Linear algebra & Convergence practice & Separation; fixed points \\
4 & Optimization & FOC/SOC; quasiconcavity & Lagrange; envelope & Farkas geometry; KKT \\
5 & Dyn.\ prog. & Framework; state \& control & Bellman equation & Backward induction; VFI idea \\
\bottomrule
\end{tabular}
\end{center}

\vspace{0.6em}
\noindent\textit{Note.} Appendices on the slides (rigorous proofs, extra examples) are for interest / self-study and are not required classroom coverage.

\end{document}
